% GENERATED FILE. Edit canonical lessons, then run npm run generate:books.
% canonical-source-hash: fnv1a64:8071553deb09f877
% canonical-lessons: MR-C160-bhavishya, MR-C160-lavkarac, MR-C160-pudhca, MR-C160-yanantar, MR-C160-akhne

\chapter{The Future, Soon, Next, From Now On, To Plan}
\label{ch:mr-the-future-soon-next-from-now-on-to-plan}
\hlchaptermodality{160}

\begin{chapteropening}
\textbf{By the end of this chapter:} \emph{I can say what will happen, and when: soon, next week, from now on.}

Complete the last lesson of chapter 160: I can say what will happen, and when: soon, next week, from now on.
\end{chapteropening}

\section[\texorpdfstring{\mr{भविष्य} (bhaviṣya)}{भविष्य (bhaviṣya)}]{\mr{भविष्य} (bhaviṣya) — the future}

\label{lesson:MR-C160-bhavishya}



\begin{quote}
Before the new one: say the Marathi for the past, then the Marathi for to happen.
\end{quote}

\subsection*{You'll want to know: \mr{भविष्य}}
\textbf{\mr{भविष्य}} — \emph{bhaviṣya} — \textquotedblleft{}the future\textquotedblright{}. Marathi's future adds endings to the verb: \textbf{\mr{मी} \mr{जाईन}} — \emph{mī jāīn} — \textquotedblleft{}I will go\textquotedblright{}; \textbf{\mr{तो} \mr{जाईल}} — \emph{to jāīl} — \textquotedblleft{}he will go\textquotedblright{}.

\begin{grammarlens}[title={What you've built}]
One more word for this chapter.
\end{grammarlens}

\subsection*{Your turn}
\begin{itemize}
\raggedright
  \item \emph{Say it:} \emph{bhaviṣya}
  \item \emph{Say it:} \emph{bhaviṣya}, once more
  \item \emph{Say it:} say \emph{bhaviṣya}
  \item \emph{Recall:} say the Marathi for the past, then the Marathi for to happen, then say \emph{bhaviṣya} again
\end{itemize}

\begin{tcolorbox}[breakable,colback=teal!4,colframe=teal!35!black,title={Before you move on}]
What does \mr{भविष्य} mean? (\textquotedblleft{}The future\textquotedblright{}.) Say it once more.
\end{tcolorbox}
\section[\texorpdfstring{\mr{लवकरच} (lavkarac)}{लवकरच (lavkarac)}]{\mr{लवकरच} (lavkarac) — soon}

\label{lesson:MR-C160-lavkarac}



\begin{quote}
Before the new one: say the Marathi for to happen, then the Marathi for the future.
\end{quote}

\subsection*{You'll want to know: \mr{लवकरच}}
\textbf{\mr{लवकरच}} — \emph{lavkarac} — \textquotedblleft{}soon\textquotedblright{}. \textbf{\mr{लवकर}}, \textquotedblleft{}early, quickly\textquotedblright{}, with the emphatic \textbf{\mr{च}}: \textbf{\mr{मी} \mr{लवकरच} \mr{येईन}.} — \emph{mī lavkarac yeīn} — \textquotedblleft{}I will come soon.\textquotedblright{}

\begin{grammarlens}[title={What you've built}]
One more word for this chapter.
\end{grammarlens}

\subsection*{Your turn}
\begin{itemize}
\raggedright
  \item \emph{Say it:} \emph{lavkarac}
  \item \emph{Say it:} \emph{lavkarac}, once more
  \item \emph{Say it:} say \emph{lavkarac}
  \item \emph{Recall:} say the Marathi for to happen, then the Marathi for the future, then say \emph{lavkarac} again
\end{itemize}

\begin{tcolorbox}[breakable,colback=teal!4,colframe=teal!35!black,title={Before you move on}]
What does \mr{लवकरच} mean? (\textquotedblleft{}Soon\textquotedblright{}.) Say it once more.
\end{tcolorbox}
\section[\texorpdfstring{\mr{पुढचा} / \mr{पुढची} / \mr{पुढचे} (puḍhcā / puḍhcī …)}{पुढचा / पुढची / पुढचे (puḍhcā / puḍhcī …)}]{\mr{पुढचा} / \mr{पुढची} / \mr{पुढचे} (puḍhcā / puḍhcī / puḍhce) — next}

\label{lesson:MR-C160-pudhca}



\begin{quote}
Before the new one: say the Marathi for the future, then the Marathi for soon.
\end{quote}

\subsection*{You'll want to know: \mr{पुढचा} / \mr{पुढची} / \mr{पुढचे}}
\textbf{\mr{पुढचा} / \mr{पुढची} / \mr{पुढचे}} — \emph{puḍhcā / puḍhcī / puḍhce} — \textquotedblleft{}next\textquotedblright{}. It agrees with its noun: \textbf{\mr{पुढच्या} \mr{आठवड्यात}} — \emph{puḍhcyā āṭhavḍyāt} — \textquotedblleft{}next week\textquotedblright{}.

\begin{grammarlens}[title={What you've built}]
One more word for this chapter.
\end{grammarlens}

\subsection*{Your turn}
\begin{itemize}
\raggedright
  \item \emph{Say it:} \emph{puḍhcā / puḍhcī / puḍhce}
  \item \emph{Say it:} \emph{puḍhcā / puḍhcī / puḍhce}, once more
  \item \emph{Say it:} say \emph{puḍhcā / puḍhcī / puḍhce}
  \item \emph{Recall:} say the Marathi for the future, then the Marathi for soon, then say \emph{puḍhcā / puḍhcī / puḍhce} again
\end{itemize}

\begin{tcolorbox}[breakable,colback=teal!4,colframe=teal!35!black,title={Before you move on}]
What does \mr{पुढचा} / \mr{पुढची} / \mr{पुढचे} mean? (\textquotedblleft{}Next\textquotedblright{}.) Say it once more.
\end{tcolorbox}
\section[\texorpdfstring{\mr{यानंतर} (yānantar)}{यानंतर (yānantar)}]{\mr{यानंतर} (yānantar) — from now on, after this}

\label{lesson:MR-C160-yanantar}



\begin{quote}
Before the new one: say the Marathi for soon, then the Marathi for next.
\end{quote}

\subsection*{You'll want to know: \mr{यानंतर}}
\textbf{\mr{यानंतर}} — \emph{yānantar} — \textquotedblleft{}from now on, after this\textquotedblright{}. It is \textbf{\mr{या}}, \textquotedblleft{}this\textquotedblright{}, with \textbf{\mr{नंतर}}, \textquotedblleft{}after\textquotedblright{}: \textbf{\mr{यानंतर} \mr{मी} \mr{रोज} \mr{चालेन}.} — \emph{yānantar mī roj cālen} — \textquotedblleft{}From now on I will walk every day.\textquotedblright{}

\begin{grammarlens}[title={What you've built}]
One more word for this chapter.
\end{grammarlens}

\subsection*{Your turn}
\begin{itemize}
\raggedright
  \item \emph{Say it:} \emph{yānantar}
  \item \emph{Say it:} \emph{yānantar}, once more
  \item \emph{Say it:} say \emph{yānantar}
  \item \emph{Recall:} say the Marathi for soon, then the Marathi for next, then say \emph{yānantar} again
\end{itemize}

\begin{tcolorbox}[breakable,colback=teal!4,colframe=teal!35!black,title={Before you move on}]
What does \mr{यानंतर} mean? (\textquotedblleft{}From now on, after this\textquotedblright{}.) Say it once more.
\end{tcolorbox}
\section[\texorpdfstring{\mr{आखणे} (ākhṇe)}{आखणे (ākhṇe)}]{\mr{आखणे} (ākhṇe) — to plan, to draw up}

\label{lesson:MR-C160-akhne}



\begin{quote}
Before the new one: say the Marathi for next, then the Marathi for from now on.
\end{quote}

\subsection*{You'll want to know: \mr{आखणे}}
\textbf{\mr{आखणे}} — \emph{ākhṇe} — \textquotedblleft{}to plan, to draw up\textquotedblright{}. Said of a plan or a line: \textbf{\mr{मी} \mr{प्रवासाची} \mr{योजना} \mr{आखीन}.} — \emph{mī pravāsācī yojanā ākhīn} — \textquotedblleft{}I will plan the trip.\textquotedblright{}

\begin{grammarlens}[title={What you've built}]
One more word for this chapter.
\end{grammarlens}

\subsection*{Your turn}
\begin{itemize}
\raggedright
  \item \emph{Say it:} \emph{ākhṇe}
  \item \emph{Say it:} \emph{ākhṇe}, once more
  \item \emph{Say it:} say \emph{ākhṇe}
  \item \emph{Recall:} say the Marathi for next, then the Marathi for from now on, then say \emph{ākhṇe} again
\end{itemize}

\begin{tcolorbox}[breakable,colback=teal!4,colframe=teal!35!black,title={Before you move on}]
What does \mr{आखणे} mean? (\textquotedblleft{}To plan, to draw up\textquotedblright{}.) Say it once more.
\end{tcolorbox}
